\section{Introduction}
\label{sec:intro}

The paradigm of deep model merging has emerged as a powerful and computationally elegant solution for multi-task learning without the prohibitive costs of joint multi-task fine-tuning \cite{wortsman2022model,ilharco2022editing,jin2022dataless}.
By directly combining the parameters of multiple task-specific expert networks that share a common pre-trained initialization, model merging creates a unified, multi-task network.
Traditional weight-space interpolation techniques—such as Task Arithmetic \cite{ilharco2022editing} or Ties-Merge \cite{yadav2023ties}—rely on static, heuristically defined, or uniform merging coefficients across all network layers.
While highly intuitive, these uniform baselines struggle when task representations conflict or when individual layers exhibit varying task-specific sensitivity.
To address this limitation, recent literature has shifted toward \emph{adaptive model merging} \cite{adamerging,regcalmerge,polymerge}.
In particular, methods like AdaMerging \cite{adamerging} employ online unsupervised Test-Time Adaptation (TTA) on incoming target streams to dynamically optimize unconstrained layer-wise coefficients by minimizing the entropy of model predictions \cite{wang2020tent}.
Under the "no-data" assumption—where no labeled validation data is available—unconstrained online TTA claims to achieve state-of-the-art multi-task performance.
Other works, such as PolyMerge \cite{polymerge}, restrict these coefficients to polynomial trajectories to reduce optimization dimensionality, yet still rely on test-time optimization over unlabeled streams.
Furthermore, because online TTA adapts on unlabeled, mixed target streams during deployment, a critical question arises: how does the system route a test sample to the correct task-specific output head?
In standard multi-task deployment, the streaming input is unlabeled, forcing the system to either minimize prediction entropy blindly across all output heads jointly (Unsupervised TTA), or rely on privileged, oracle task-routing labels at inference time to direct gradients and predictions to the active task head (a setting we call the "privilege trap").
More recently, critics of online TTA have proposed Offline Few-Shot Validation Tuning (OFS-Tune), demonstrating that in practical scenarios, a tiny labeled validation set (e.g., $M=10$ samples per task) can be leveraged offline to optimize stable layer-wise trajectories, eliminating test-time computational overhead and latency while naturally bypassing this task-routing dilemma entirely with respect to weight-adaptation and gradient backpropagation (since the merged model parameters are static, preventing representation collapse under joint optimization).
However, we explicitly qualify that while OFS-Tune completely eliminates routing requirements for parameter adaptation, any multi-head merged model deployed on interleaved streams still requires a routing mechanism at inference time to select the correct task-specific output head for final prediction, which remains a shared challenge for all multi-task architectures.
We emphasize from the outset that online TTA and offline OFS-Tune operate under fundamentally distinct data-access assumptions and deployment settings, representing a key trade-off for practitioners.
Online TTA is designed for strictly unsupervised, zero-shot environments where labeled training or validation data is entirely unavailable or cannot be stored due to user-privacy regulations.
In these settings, online TTA offers maximum operational flexibility by adapting dynamically on live, unlabeled test streams without requiring prior developer calibration.
Conversely, offline OFS-Tune assumes access to a minimal labeled calibration set (e.g., 10 samples per task) prior to deployment.
This mild assumption enables OFS-Tune to directly optimize supervised classification loss offline, acting as a powerful analytical filter that rejects transductive noise and completely eliminates test-time compute, backpropagation latency, and edge-device energy consumption.
By contrasting these two paradigms across varying task conflicts, we provide practitioners with a comprehensive map of when each paradigm is operationally safe and effective.
We examine this rapid research cycle with constructive skepticism.
We identify a major, un-reported \textbf{confounding variable and hidden assumption} in this entire body of work: \textbf{the standard evaluation protocol is overfitted to a single, arbitrary task suite}.
Indeed, almost all existing adaptive model-merging publications validate their proposed algorithms on exactly one combination of four datasets: MNIST, FashionMNIST, CIFAR-10, and SVHN.
This specific combination is highly heterogeneous and arbitrary, spanning simple grayscale digits, grayscale fashion items, street numbers, and complex natural RGB objects.
We ask: \textbf{Is this single evaluation protocol robust?
Or does the relative ranking of model-merging methods suffer from a severe "Task Suite Bias" that masks fundamental, catastrophic failures under different task relationships?}

To answer this, we present \textbf{SuiteMerge}, a systematic and rigorous methodological audit of modern adaptive model-merging paradigms.
We systematically partition the standard four-dataset pool into five distinct multi-task evaluation suites designed along axes of domain distance and representational conflict:
\begin{itemize}
    \item \textbf{Suite A (Highly Homogeneous - Low Conflict):} MNIST + FashionMNIST.
Representational overlap friction is extremely low.
    \item \textbf{Suite B (Highly Heterogeneous - High Conflict):} CIFAR-10 + SVHN.
Severe representational clashes occur in shared network subspaces.
    \item \textbf{Suite C (Cross-Domain Digits):} MNIST + SVHN.
Digital classification featuring a massive domain shift from clean grayscale to cluttered, natural street scenes.
    \item \textbf{Suite D (Cross-Domain Objects):} FashionMNIST + CIFAR-10.
Grayscale fashion items contrasted against RGB natural objects.
    \item \textbf{Suite E (Full 4-Task Suite - Control):} MNIST + FashionMNIST + CIFAR-10 + SVHN.
The standard control benchmark used in previous work.
\end{itemize}

By evaluating Uniform merging, online AdaMerging, online PolyMerge, and offline OFS-Tune under realistic stream noise and batch bias, our rigorous audit exposes three critical methodological revelations:
First, the relative ranking of methods is highly sensitive to the chosen task suite.
On highly homogeneous tasks (Suite A), naive Uniform merging is extremely competitive, and TTA functions reasonably well.
Second, on highly heterogeneous, high-conflict suites (Suite B and Suite D), unconstrained online TTA (AdaMerging) overfits to local transductive stream noise.
By optimizing high-dimensional unconstrained layer-wise coefficients (e.g., 48 parameters for a 12-layer ViT), the unsupervised entropy-minimization objective over-parameterizes on stream noise.
While this results in a mild performance lag behind regularized and offline methods in simulation (where it still outperforms Uniform), it causes catastrophic representation collapse below the static Uniform baseline in actual physical weight-space networks.
Third, offline OFS-Tune using continuous low-degree polynomials behaves as a powerful analytical low-pass filter.
By constraining the optimization space offline on $M=10$ labeled samples, it completely rejects validation sample noise and transductive domain shifts.
We propose both a linear profile ($d=1$) and a quadratic profile ($d=2$) as primary configurations.
In simulation, OFS-Tune consistently outperforms unconstrained online TTA (AdaMerging) by up to \textbf{5.33\%} in accuracy on high-conflict suites.
With a linear profile ($d=1$), OFS-Tune matches the complex online PolyMerge ($d=2$) to within $1.02\%$ across all suites, while our quadratic configuration ($d=2$) matches PolyMerge to within $0.26\%$ on homogeneous/mixed suites and successfully outperforms it in the high-conflict Suite B (achieving \textbf{68.62\%} vs.
68.51\% average accuracy) while reducing standard deviation.
Crucially, in real physical weight-space networks, OFS-Tune successfully outperforms online PolyMerge by up to \textbf{3.70\%} and online AdaMerging by up to \textbf{4.20\%} because it completely eliminates transductive stream-level overfitting and avoids the privileged task-routing assumptions that online methods require to function without collapsing.
This is achieved whilst requiring \textbf{zero test-time compute and zero latency overhead} at deployment.
Why has this catastrophic failure of online TTA been obscured in previous publications?
Our audit reveals a subtle mathematical mask: in the full 4-task control suite (Suite E), the massive structural distance of Uniform merging from any individual task's optimal coefficient profile creates a large initial accuracy degradation.
This large gap makes the over-parameterized online TTA optimizer appear successful simply because it moves coefficients away from this poor uniform starting point, masking the fact that it suffers from localized representation collapse on high-conflict task components.
In summary, our key contributions are:
\begin{enumerate}
    \item We expose the \textbf{Task Suite Bias} in the model-merging literature, demonstrating that a single multi-task suite is a dangerous confounding variable that leads to false conclusions regarding algorithm generalizability.
    \item We mathematically formulate and empirically demonstrate \textbf{transductive overfitting} in unconstrained online TTA under realistic, correlated stream noise, proving that unsupervised test-time optimization collapses when representational conflicts are high.
    \item We establish that offline polynomial validation tuning (\textbf{OFS-Tune}) acts as a robust analytical regularizer, providing superior generalization across all task relationships with zero deployment overhead.
    \item We present a comprehensive, multi-seed comparative benchmark across five distinct suites, advocating for a transition toward multi-suite validation standards in model-merging research.
\end{enumerate}
